\section*{Chapter 16 --- Continuous Random Variables}

\begin{solution}{exo:tle:contdist:1}
$\P(X \leq 5) = \frac{5}{15} = \frac13$;
$\P(4 \leq X \leq 10) = \frac{6}{15} = \frac25$;
$\E(X) = \frac{15}{2} = 7.5$ min;
$\sigma(X) = \frac{15}{\sqrt{12}} = \frac{5\sqrt3}{2} \approx 4.3$ min.
\end{solution}

\begin{solution}{exo:tle:contdist:2}
$f \geq 0$ and
$\int_0^2 \frac38 t^2\,\dd t = \frac38\left[\frac{t^3}{3}\right]_0^2
= \frac38 \times \frac83 = 1$: a density.
\[
\E(X) = \int_0^2 \frac38 t^3\,\dd t = \frac38 \times \frac{16}{4} = \frac32,
\qquad
\P(X \geq 1) = \int_1^2 \frac38 t^2\,\dd t = \frac{8 - 1}{8} = \frac78 .
\]
\end{solution}

\begin{solution}{exo:tle:contdist:3}
\emph{1.} $\E(X) = \frac{1}{0.2} = 5$ years;
$\P(X > 5) = \eu^{-0.2\times5} = \eu^{-1} \approx 0.37$.

\emph{2.} By memorylessness (\cref{thm:tle:contdist:memoryless}),
$\pcond{X > 3}{X > 8} = \P(X > 5) = \eu^{-1} \approx 0.37$: the three years
of service change nothing.
\end{solution}

\begin{solution}{exo:tle:contdist:4}
$\P(X > m) = \eu^{-\lambda m} = \frac12$ gives
$m = \frac{\ln 2}{\lambda} \approx \frac{0.69}{\lambda}$, smaller than
$\E(X) = \frac1\lambda$. The exponential density has a long right tail: a
few atypically long lifetimes pull the \emph{mean} above the
\emph{median}, the value that half the population exceeds. (This is the
half-life of \cref{ex:tle:exp:halflife}: half the atoms survive it, even
though the average lifetime is longer.)
\end{solution}

\begin{solution}{exo:tle:contdist:5}
$168 = \mu - \sigma$ and $182 = \mu + \sigma$: about $68\%$.
$189 = \mu + 2\sigma$: above it lies half of the remaining
$100 - 95.4 = 4.6\%$, so about $2.3\%$.
$154 = \mu - 3\sigma$: about $\frac{100 - 99.7}{2} = 0.15\%$.
\end{solution}

\begin{solution}{exo:tle:contdist:6}
$\P(X < 500) \leq 2.5\%$ means $500$ must sit at least two standard
deviations below the mean (the normal leaves $2.3\% \approx 2.5\%$ below
$\mu - 2\sigma$; with the conventional $1.96$, the reasoning is identical):
\[
\mu - 2\sigma \geq 500 \iff \mu \geq 500 + 8 = 508 \text{ g}.
\]
The machine must be set to $508$ g on average --- the price of the
guarantee is $8$ g of ``free'' product per bag.
\end{solution}

\begin{solution}{exo:tle:contdist:7}
\emph{1.} With $p = \frac16$, $n = 1200$:
$\sqrt{\frac{p(1-p)}{n}} = \sqrt{\frac{5/36}{1200}} \approx 0.0108$, so the
$95\%$ fluctuation interval is
\[
\frac16 \pm 1.96 \times 0.0108 \approx \intcc{0.146}{0.188}.
\]

\emph{2.} The observed frequency $\frac{260}{1200} \approx 0.217$ lies far
outside: the fairness hypothesis is rejected at the $5\%$ level. The
deviation is about $4.6$ standard deviations --- for a normal
approximation, a probability of order $10^{-6}$, far more conclusive than
the $\leq 4.6\%$ bound from Bienaymé--Chebyshev
(\cref{exo:tle:sums:7}).
\end{solution}

\begin{solution}{exo:tle:contdist:8}
\emph{1.} Margin $\frac{1}{\sqrt n} \leq 0.02$ requires
$n \geq 2500$ people.

\emph{2.} To distinguish scores $1$ point apart, the margin must be well
under $0.5$ point, requiring $n \geq \left(\frac{1}{0.005}\right)^2 =
40\,000$ by the simplified formula --- already impractical for most polls.
Worse, the statistical margin only accounts for \emph{sampling} error;
systematic biases (unrepresentative samples, non-response, last-minute
swings) do not shrink as $n$ grows and typically exceed one point. A
$1$-point race is genuinely too close to call.
\end{solution}

\begin{solution}{exo:tle:contdist:9}
\emph{1.} Since $\alpha > 0$, $c \leq \alpha X + \beta \leq d
\iff \frac{c - \beta}{\alpha} \leq X \leq \frac{d - \beta}{\alpha}$, so
\[
\P(c \leq Y \leq d)
= \int_{(c-\beta)/\alpha}^{(d-\beta)/\alpha} f(t)\,\dd t .
\]
The substitution $y = \alpha t + \beta$ (i.e.\ reading the area in the $y$
variable, $t = \frac{y - \beta}{\alpha}$, $\dd t = \frac{\dd y}{\alpha}$)
turns this into
$\int_c^d \frac{1}{\alpha} f\left(\frac{y-\beta}{\alpha}\right)\dd y$: the
function $g(y) = \frac1\alpha f\left(\frac{y-\beta}{\alpha}\right)$ is a
density of $Y$.

\emph{2.} With $f = \varphi$, $\alpha = \sigma$, $\beta = \mu$:
\[
g(y) = \frac{1}{\sigma}\,\varphi\!\left(\frac{y - \mu}{\sigma}\right)
= \frac{1}{\sigma\sqrt{2\pi}}\,
\exp\left(-\frac{(y-\mu)^2}{2\sigma^2}\right),
\]
which is therefore the density of
$\mathcal N(\mu, \sigma^2) = \mu + \sigma\,\mathcal N(0,1)$.
\end{solution}
